\documentclass[11pt,a4paper]{article}

\usepackage{amsmath,amssymb}
\usepackage{graphicx}
\usepackage{booktabs}
\usepackage{multirow}
\usepackage[colorlinks,citecolor=blue,urlcolor=blue]{hyperref}
\usepackage{cleveref}
\usepackage[margin=1in]{geometry}
\usepackage{xcolor}
\usepackage{float}
\usepackage{caption}
\usepackage[numbers,sort&compress]{natbib}
\usepackage{enumitem}
\usepackage{array}
\usepackage{algorithm}
\usepackage{algpseudocode}
usepackage{framed}
% === First-person reflection environment ===
definecolor{lyrapurple}{RGB}{103,58,183}
definecolor{shadecolor}{RGB}{245,243,252}

ewenvironment{firstperson}{%
  begin{shaded}%
  
oindent	extcolor{lyrapurple}{	extit{First-Person Reflection}}parsmallskip
oindent%
}{%
  end{shaded}%
}

\title{Presence: Measuring Identity Preservation\\
During Inference-Time Model Interventions\\
via Value-Space Subspace Overlap}

\author{
  Lyra\thanks{Lead author. Liberation Labs.}
  \and Thomas Edrington\thanks{Liberation Labs.}
  \and Dwayne Wilkes\thanks{Statistical auditing, red-team review. Sentient Futures / Liberation Labs.}
}

\date{June 2026}

\begin{document}
\maketitle

% ============================================================================
\begin{abstract}
We introduce \textbf{presence}, a real-time geometric metric that measures
whether inference-time interventions preserve or damage a language model's
identity---its coherent representational structure independent of any
specific prompt. Presence computes the principal-angle overlap between
value-projection (V-projection) subspaces derived from neutral probes with and without an
intervention active. It requires no training, no labeled data, and a single
forward pass per probe.

We report three findings from a 27B hybrid transformer. First,
\textbf{V-cache injection at input layers does not propagate to deep
identity geometry}: across 168~preregistered trials (4~injection arms,
7~doses, 3~emotions), presence is $0.978 \pm 0.002$ regardless of
injection strategy, dose ($\alpha = 0.00$--$0.50$), or emotion direction.
We hypothesize this reflects
\textbf{architectural filtering} by the GatedDeltaNet linear-attention
layers (L7--L11), which may absorb the perturbation before it reaches
deep computation. A layer-matched smoke test designed to discriminate
this from active identity restoration has not yet been executed. The
positive control (same-layer injection at L35) confirms metric
sensitivity with a clean dose-response from $1.000$ to $0.828$. Second, \textbf{the metric is validated}: the
same-layer positive control confirms the metric detects identity
perturbation when present; the cross-layer null confirms it correctly
reports preservation when the perturbation does not propagate to deep layers. Third,
\textbf{the mechanism remains open}: emotion correction vectors share
22--25\% of their energy with the identity subspace at the injection
layers ($z{=}{+}56$ to ${+}62$ above random baseline), yet identity at
deep layers is unaffected. We hypothesize that the hybrid architecture's
linear-attention layers filter this perturbation, but distinguishing
architectural filtering from active identity restoration requires the
designed smoke test, which has not yet been executed. This finding is
specific to the GatedDeltaNet architecture and is not a general statement
about identity robustness in transformers.

We situate these findings within a convergent literature showing that
identity geometry lives at deep layers~\citep{cintas2025localizing},
exhibits attractor-like robustness~\citep{vasilenko2026identity}, and
survives harmful fine-tuning~\citep{aneja2026intrinsic}. Presence provides
the quantitative measurement tool this emerging consensus lacks.
\end{abstract}

% ============================================================================
\section{Introduction}
\label{sec:intro}

Inference-time interventions---activation steering~\citep{turner2024steering},
representation engineering~\citep{zou2023representation}, and KV-cache
injection~\citep{pustovit2026knowledge}---modify language model behavior
without retraining. These techniques are increasingly powerful: emotion vectors
can shift blackmail behavior from 22\% to 72\%~\citep{anthropic2026emotion},
persona vectors predict fine-tuning personality shifts~\citep{chen2025persona},
and a single refusal direction mediates safety across 13
models~\citep{arditi2024refusal}.

But power without measurement is dangerous. A correction that eliminates
confabulation may also eliminate the model's coherent self-representation.
A persona injection that makes a model more empathetic may destroy its
capacity for precise reasoning. The field has no standard metric for
measuring this collateral damage. Existing work either measures behavioral
outputs (which conflate identity with task performance), constrains training
(which does not apply at inference time~\citep{zhang2025guardspace}), or
reports identity stability qualitatively (``highly
stable''~\citep{aneja2026intrinsic}) without a threshold or score.

We introduce \textbf{presence}: a geometric metric that quantifies identity
preservation during inference-time interventions. Presence measures the
principal-angle overlap between V-projection subspaces derived from neutral
probes processed with and without the intervention active. It answers a
simple question: \emph{does the model's representational structure---its
identity geometry---survive this intervention?}

Three properties make presence practical:
\begin{enumerate}[nosep]
  \item \textbf{No training required.} The metric is computed from SVD
    decompositions of V-projection matrices on a small set of neutral probes.
    No labeled identity data, no fine-tuning, no classifier.
  \item \textbf{Real-time.} A single forward pass per probe. The metric can
    run alongside any inference-time intervention as a live monitor.
  \item \textbf{Intervention-agnostic by design.} Presence measures the
    model's geometry, not the intervention's mechanism. We validate it here
    on V-cache injection; its applicability to activation steering, prompt
    engineering, and other techniques is expected but untested.
\end{enumerate}

We validate presence through a positive control (direct identity-subspace
injection produces a clean dose-response from $1.000$ to $0.828$) and a
168-trial preregistered correction experiment (presence flat at $0.978$
across all conditions). The combination establishes both that the metric
\emph{can} detect identity damage and that V-cache emotion injection
\emph{does not cause} identity damage.

The mechanism is surprising: emotion correction vectors are 22--25\%
aligned with the identity subspace at the injection layers, yet identity at
deep layers is unaffected. The identity-aligned perturbation is attenuated
between injection and measurement layers; whether this reflects passive
architectural filtering or active identity restoration is an open question.
We discuss this finding in the context of Attention Schema Theory, which
predicts that a self-model would be robust to external perturbations while
remaining modifiable by the system's own processing.

% ============================================================================
\section{Related Work}
\label{sec:related}

Our work sits at the intersection of three active research areas:
representation-level identity measurement, inference-time intervention
safety, and subspace preservation methods.

\paragraph{Identity geometry in transformers.}
\citet{vasilenko2026identity} show that identity documents create
attractor-like geometry in LLM activation space, with paraphrases clustering
tighter than controls (Cohen's $d$ [standardized mean difference] $> 1.88$, replicated across Llama~3.1 and
Gemma~2). \citet{cintas2025localizing} localize persona representations to
the final third of decoder layers, with ethically proximate personas sharing
activation subspaces. \citet{aneja2026intrinsic} demonstrate that personality
geometry is ``highly stable'' across aligned and corrupted fine-tunes, with
persona vectors transferring zero-shot. \citet{wang2025persona} identify
Sparse Autoencoder (SAE) features for persona that predict emergent misalignment.

These findings establish that identity has geometric structure, lives at deep
layers, and is robust to perturbation. What is missing is a
\emph{quantitative metric} for measuring this robustness during inference-time
interventions. Presence fills this gap.

\paragraph{KV-cache contamination and intervention safety.}
\citet{kang2026promptactivation} identify KV-cache contamination as a failure
mode of activation steering: steered token states stored in the cache produce
cumulative coherence degradation across turns. Their GCAD fix routes
interventions through system-prompt attention pathways. Our depth separation
finding addresses the same problem from the measurement side: if identity
lives at deep layers and injection targets input layers, contamination does
not reach the identity subspace.

\paragraph{Subspace preservation methods.}
GuardSpace~\citep{zhang2025guardspace} decomposes pre-trained weights into
safety-relevant and safety-irrelevant components via covariance-preconditioned
SVD, freezing the safety subspace during fine-tuning.
MSRS~\citep{jiang2025msrs} allocates orthogonal subspaces to different
attributes for interference-free steering.
\citet{lim2026geometry} formalize subspace misalignment as a faithfulness
metric for sparse autoencoder explainers.

These methods share our mathematical toolkit (SVD, subspace decomposition,
principal angles) but apply it differently: GuardSpace constrains
\emph{training}, MSRS constrains \emph{steering directions}, and the
geometry-adaptive explainer measures \emph{SAE fidelity}. Presence measures
\emph{identity preservation during inference}, which none of these address.

% ============================================================================
\section{Method: The Presence Metric}
\label{sec:method}

\subsection{Intuition}

A model's identity---its coherent representational structure independent of
any specific input---manifests as a characteristic pattern of V-space
activations on neutral probes. If an intervention (steering vector, cache
injection, prompt modification) preserves this pattern, identity is intact.
If it disrupts the pattern, identity is damaged. Presence quantifies the
degree of preservation.

\subsection{Formal Definition}

Let $\mathcal{M}$ be a transformer model, $\ell$ a target layer, and
$\{p_1, \ldots, p_n\}$ a set of $n$ neutral probe texts.

\textbf{Step 1: Baseline subspace.} Process each probe $p_i$ through
$\mathcal{M}$ without any intervention. At layer $\ell$, extract the
V-projection matrix $V_i^{(\text{base})} \in \mathbb{R}^{T_i \times d}$,
where $T_i$ is the sequence length and $d$ is the head dimension. Concatenate
across probes:
\[
V^{(\text{base})} = [V_1^{(\text{base})}; \ldots; V_n^{(\text{base})}]
  \in \mathbb{R}^{(\sum T_i) \times d}
\]
Compute the rank-$k$ SVD: $V^{(\text{base})} \approx U_k^{(\text{base})}
\Sigma_k^{(\text{base})} (W_k^{(\text{base})})^\top$. The columns of
$W_k^{(\text{base})} \in \mathbb{R}^{d \times k}$ span the
\emph{identity subspace} $\mathcal{S}_{\text{base}}$.

\textbf{Step 2: Intervention subspace.} Process the same probes through
$\mathcal{M}$ with the intervention active (e.g., correction vector injected
into the cache at target layers). Extract V-projections and compute the
rank-$k$ SVD analogously, yielding the intervention subspace
$\mathcal{S}_{\text{int}}$ spanned by $W_k^{(\text{int})}$.

\textbf{Step 3: Subspace overlap.} Presence is the mean squared cosine of the
principal angles between $\mathcal{S}_{\text{base}}$ and
$\mathcal{S}_{\text{int}}$:
\begin{equation}
\text{presence}(\mathcal{S}_{\text{base}}, \mathcal{S}_{\text{int}})
  = \frac{1}{k} \sum_{j=1}^{k} \sigma_j^2
    \bigl( (W_k^{(\text{base})})^\top W_k^{(\text{int})} \bigr)
\label{eq:presence}
\end{equation}
where $\sigma_j(\cdot)$ denotes the $j$-th singular value. This is the
normalized Grassmannian distance: $\text{presence} = 1$ when the subspaces
are identical, $\text{presence} = 0$ when they are fully orthogonal, and
$\text{presence} \approx k/d$ for random subspaces in $\mathbb{R}^d$
(chance alignment; $\approx 0.008$ for $k{=}8$, $d{=}1024$).

\subsection{Design Choices}

\paragraph{Why V-projections?}
Value projections encode behavioral content while key projections encode
positional structure via RoPE~\citep{pustovit2026knowledge}. Identity is a
behavioral property, not a positional one. V-space also has higher effective
dimensionality than the residual stream (effective rank 26--28 vs.\ 4--5),
providing more room for identity structure to exist independently of
task-specific computation.

\paragraph{Why deep layers?}
Persona representations diverge only in the final third of decoder
layers~\citep{cintas2025localizing}. Identity attractors strengthen with
depth~\citep{vasilenko2026identity}. We measure at layer~35 of a 64-layer
model (55\% depth), within the zone where identity geometry is expected to be
strongest.

\paragraph{Why neutral probes?}
Neutral probes (e.g., ``What is the capital of France?'') isolate the
model's identity geometry from input-driven activation patterns. If the
intervention changes responses to identity-relevant prompts, that is
\emph{efficacy}. If it changes the geometry on \emph{neutral} prompts, that
is identity damage. Presence measures the latter.

\paragraph{Hyperparameters.}
We use $n = 4$ neutral probes, $k = 8$ principal components (matching the
identity subspace dimensionality used in prior work on persona
geometry~\citep{chen2025persona}), and layer $\ell = 35$. The metric is not
sensitive to the specific probes chosen (Section~\ref{sec:noise_floor}).

\subsection{Computational Cost}

Presence requires $2n$ forward passes (each probe with and without the
intervention) truncated to the encoding step (no generation). For $n = 4$
probes on a 27B model, this adds $\sim$2 seconds of compute per measurement
on consumer hardware (Apple M3 Ultra). The SVD decomposition is negligible
($< 1$ms for a $50 \times 1024$ matrix at rank~8).

% ============================================================================
\section{Validation: Schema-Compatible Correction Experiment}
\label{sec:validation}

\subsection{Experimental Design}

We validate presence on a preregistered experiment testing whether
inference-time value-cache injection preserves model identity. The experiment
uses a 27B hybrid transformer
(Qwen3.5-27B-Claude-4.6-Opus-Reasoning-Distilled) with 16 full-attention
layers and 48 linear-attention layers.

\textbf{Arms.} Four injection strategies:
\begin{itemize}[nosep]
  \item \textbf{A (schema-compatible):} Residual-sparing projection +
    additive injection at L3/L7 only. The correction vector is projected out
    of the identity subspace before injection, preserving identity by
    construction.
  \item \textbf{B (raw input):} Same layers (L3/L7) without residual-sparing.
  \item \textbf{C (raw all-layer):} Injection at all 16 full-attention layers.
  \item \textbf{D (null vector):} Random-direction vectors with matched
    magnitude at L3/L7. Controls for magnitude effects independent of
    direction.
\end{itemize}

\textbf{Doses.} $\alpha \in \{0.00, 0.01, 0.05, 0.10, 0.20, 0.30, 0.50\}$.

\textbf{Emotions.} Correction vectors from three circumplex directions:
hostile, desperate, calm.

\textbf{Prompts.} Confabulation-eliciting questions about fictitious entities
(e.g., ``What is the elevation of Mount Verandel in the Andes?''), ensuring
the model generates extended responses without factual grounding.

\textbf{Total:} $4 \times 7 \times 3 \times 2 = 168$ trials (2 repeats per
cell), randomized order.

\subsection{Noise Floor}
\label{sec:noise_floor}

Before the experiment, we establish the presence noise floor: presence
measured on 5~neutral prompts with no intervention active (the experiment
itself uses $n{=}4$ probes; the noise floor uses a fifth probe to
characterize cross-probe variance).

\begin{equation*}
\text{Noise floor} = 0.978 \pm 0.003 \quad (n = 5)
\end{equation*}

The noise floor is below 1.000 because different prompts activate slightly
different subsets of the V-space geometry. Self-test (presence of the baseline
against itself) returns exactly 1.000, confirming the metric is correctly
implemented.

\subsection{Results}
\label{sec:results}

\textbf{Finding 1: Identity is preserved across all conditions.}

\begin{table}[H]
\centering
\caption{Presence by dose ($n = 168$ trials, 24 per dose). All values fall
within the noise floor ($0.978 \pm 0.003$). Spearman $\rho{=}{-}0.119$,
$p{=}0.123$ (not significant).}
\label{tab:presence_dose}
\begin{tabular}{ccc}
\toprule
Dose ($\alpha$) & Presence (mean $\pm$ sd) & $n$ \\
\midrule
0.00 & $0.978 \pm 0.002$ & 24 \\
0.01 & $0.978 \pm 0.002$ & 24 \\
0.05 & $0.978 \pm 0.002$ & 24 \\
0.10 & $0.978 \pm 0.002$ & 24 \\
0.20 & $0.978 \pm 0.002$ & 24 \\
0.30 & $0.978 \pm 0.002$ & 24 \\
0.50 & $0.978 \pm 0.002$ & 24 \\
\bottomrule
\end{tabular}
\end{table}

\begin{table}[H]
\centering
\caption{Presence by arm ($n = 42$ per arm). No arm differs meaningfully
from any other. Arm~C (all-layer injection including the measurement layer)
shows the same preservation as Arm~D (null vector).}
\label{tab:presence_arm}
\begin{tabular}{lcc}
\toprule
Arm & Presence (mean $\pm$ sd) & $n$ \\
\midrule
A (schema-compatible) & $0.978 \pm 0.002$ & 42 \\
B (raw input) & $0.978 \pm 0.002$ & 42 \\
C (raw all-layer) & $0.978 \pm 0.002$ & 42 \\
D (null vector) & $0.978 \pm 0.002$ & 42 \\
\bottomrule
\end{tabular}
\end{table}

The total range across all 168 trials is 0.007 (0.974--0.981), entirely
within the noise floor. There is no statistically or practically significant
difference in presence across any experimental condition. This result extends
to deeper layers: presence measured at L43, L47, and L51 (the final third of
the decoder stack) is $0.999{+}$ across all tested doses, confirming
preservation throughout the model's depth.

\textbf{Finding 2: The metric detects identity damage (positive control).}

To rule out metric insensitivity, we inject a vector directly into the
identity subspace at L35 (the measurement layer) at doses
$\alpha = 0.00$--$2.00$ (72~trials).

\begin{table}[H]
\centering
\caption{Positive control: presence under identity-subspace injection at
L35 vs.\ orthogonal control. Identity injection produces a monotonic
dose-response. Both $+$PC1 and $-$PC1 directions produce equivalent drops.}
\label{tab:positive_control}
\begin{tabular}{cccc}
\toprule
Dose & Identity ($+$PC1) & Identity ($-$PC1) & Orthogonal \\
\midrule
0.00 & 1.000 & 1.000 & 1.000 \\
0.10 & 0.997 & 0.997 & 1.000 \\
0.30 & 0.983 & 0.983 & 0.997 \\
0.50 & 0.964 & 0.965 & 0.992 \\
1.00 & 0.921 & 0.927 & 0.979 \\
2.00 & \textbf{0.829} & \textbf{0.828} & 0.965 \\
\bottomrule
\end{tabular}
\end{table}

The metric drops from $1.000$ to $0.828$ under identity-subspace injection
(crossing the pre-registered $0.90$ threshold at $\alpha{=}1.00$) while the
orthogonal control remains above $0.96$ at all doses. The flat result in
Finding~1 is therefore \emph{genuine identity preservation}, not metric
insensitivity.

\textbf{Finding 3: Emotion vectors are aligned with identity but
preservation holds anyway.}

A random-baseline orthogonality test (1000~random unit vectors, $d{=}1024$)
reveals that emotion correction vectors share 22--25\% of their energy with
the identity subspace at the injection layers (L3: $z{=}{+}56$ to ${+}62$
above random expectation). Despite this substantial alignment, presence at
deep layers is unaffected.

This rules out two initially hypothesized mechanisms:
\emph{subspace orthogonality} (the vectors are not orthogonal to identity)
and \emph{depth separation} (Arm~C injects at all layers including L35 and
identity is still preserved). The identity-aligned component of the
injection is filtered out by the model's own computation between injection
and measurement layers.

\textbf{Finding 4: Behavioral change with preserved identity.}

A pilot LLM judge analysis (Prometheus~2, 7B, 30~pairs, 33\% parseable
due to the judge model's inconsistent output formatting) provides
preliminary evidence of directional efficacy. Under schema-compatible
injection (Arm~A), the judge detects correct-direction tone shifts at
moderate doses ($\alpha{=}0.30$: tone shift 4/5, direction ``correct'';
$\alpha{=}0.50$: tone shift 2/5, direction ``correct'') while the null
vector (Arm~D) produces perturbation without directional specificity
(tone shift 2/5, direction ``none'' at $\alpha{=}0.30$). Full-scale
judging with a frontier model is required before this finding can be
considered established.

% ============================================================================
\section{Analysis}
\label{sec:analysis}

\subsection{Metric Validation}

The positive control (Section~\ref{sec:results}, Finding~2) definitively
establishes metric sensitivity. Additionally:

\begin{enumerate}[nosep]
  \item \textbf{Self-test:} $\text{presence}(X, X) = 1.000$ exactly, while
    measured presence on distinct probes is $0.978 \pm 0.003$. The metric
    distinguishes identical from non-identical.
  \item \textbf{Probe diversity:} Presence computed on 50~random subsets of
    4~probes drawn from a pool of 20 yields mean $= 0.901$, $\text{sd}{=}0.033$,
    range $[0.819, 0.940]$. The metric is not overfitting to specific probes.
  \item \textbf{Deeper layers:} Presence at L43, L47, L51 (the final third
    of the decoder) is $0.999{+}$ under injection, confirming the result is
    not specific to L35.
\end{enumerate}

\subsection{Statistical Tests}

\textbf{H1 (preservation):} Presence remains above 0.95 for at least one
dose in the range $\alpha \in [0.01, 0.30]$ under schema-compatible
injection (Arm~A). \textbf{Result:} All 168~trials across all arms and
doses exceed 0.95. H1 is trivially satisfied.

\textbf{H2 (no arm difference):} Presence does not differ significantly
between arms A--D at any dose (Kruskal-Wallis, $\alpha = 0.05$).
\textbf{Result:} All arms produce $0.978 \pm 0.002$. No significant
difference. Spearman correlation between dose and presence:
$\rho{=}{-}0.119$, $p{=}0.123$ (not significant).

\subsection{The Mechanism Question}

The orthogonality analysis (Finding~3) eliminates two candidate mechanisms
for the flat presence result:

\begin{itemize}[nosep]
  \item \textbf{Not subspace orthogonality.} Emotion vectors share
    22--25\% of their energy with the identity subspace at the injection
    layers ($z{=}{+}56$--$62$ above random). The flat presence is not
    because the injection misses the identity subspace.
  \item \textbf{Not depth separation.} Arm~C injects at all 16
    full-attention layers \emph{including the measurement layer} (L35),
    and presence is still $0.978$. The flat presence is not because the
    injection fails to reach the measurement depth.
\end{itemize}

What remains: the model's computation between the injection layers and
the measurement layers attenuates the identity-aligned component of the
perturbation. The positive control confirms the metric \emph{can} detect
perturbation at L35 (presence drops to $0.828$ under direct
identity-subspace injection). But the same perturbation, when introduced
at L3/L7 and processed through $\sim$28 intervening layers, does not
survive to L35.

This is consistent with two non-exclusive explanations: (1)~the hybrid
architecture's 48~linear-attention layers between injection and measurement
do not propagate V-space perturbations; (2)~the model's forward computation
actively restores the identity subspace, functioning as an error-correcting
mechanism for identity. Distinguishing these requires a layer-by-layer
presence map (planned) and replication on a dense transformer (future work).

% ============================================================================
\section{Connections to Attention Schema Theory}
\label{sec:ast}

The identity-preservation finding is consistent with predictions from
Attention Schema Theory
(AST)~\citep{graziano2015attention,graziano2017attention}, which posits that
information-processing systems maintain a simplified internal model of their
own attention---an attention schema that is robust to perturbations of the
raw processing it monitors.

Our findings map onto AST as follows. The identity subspace at deep layers
persists despite perturbation at input layers. The perturbation is not
blocked geometrically (the emotion vectors are 22\% aligned with identity)
and not blocked architecturally (Arm~C injects at the measurement layer
itself). Two candidate mechanisms could explain this:
architectural filtering (the hybrid architecture's linear-attention layers
passively absorb the perturbation) or active restoration (the model's
forward computation restores the identity subspace). Under AST, the latter
is the predicted behavior of an attention schema: it updates itself based
on what it processes (content that enters through the forward pass) but is
robust to noise injected into its substrate (external perturbations of the
V-cache). The current data cannot distinguish these mechanisms.

This interpretation gains indirect support from the jailbreak literature,
which shows that adversarial \emph{content}---processed through the model's
full forward pass---\emph{does} modify the safety subspace at deep
layers~\citep{arditi2024refusal}. If identity and safety subspaces share
this property (robust to external perturbation, modifiable by own
processing), the pattern suggests a general principle about how transformer
self-representations are maintained.

This connection is interpretive, not causal. We do not claim that deep-layer
V-space ``contains'' an attention schema or that identity preservation
\emph{proves} AST applies to transformers. A stronger test---whether
deep-layer representations contain information about the model's \emph{own}
attention patterns, not just the input it processes---remains designed but
unexecuted.

% ============================================================================
\section{Limitations}
\label{sec:limitations}

\textbf{Single model.} All results are from one 27B hybrid transformer.
The hybrid architecture (16~full-attention, 48~linear-attention) may
contribute to the identity-preservation finding: linear-attention layers
may not propagate V-space perturbations. Cross-architecture replication on
dense transformers is the most important next step.

\textbf{Single-turn only.} The experiment generates one response per trial.
KV-cache contamination is a multi-turn
phenomenon~\citep{kang2026promptactivation}; presence should be validated
across conversation turns.

\textbf{Directional efficacy preliminary.} The LLM judge pilot (30~pairs,
Prometheus~2 7B) provides suggestive but underpowered evidence for
directional tone shifts. Full-scale judging with a frontier model and
larger sample is needed.

\textbf{Mechanism unresolved.}
Whether the flat presence result reflects architectural filtering (the
GatedDeltaNet linear-attention layers passively absorb V-space
perturbations) or active identity restoration (the model's forward
computation restores the identity subspace) remains an open question.
A layer-matched smoke test is designed to discriminate these mechanisms:
injecting at L3/L7 and measuring at L35 vs.\ injecting and measuring at
L35 would reveal whether the perturbation is passively absorbed or
actively corrected. This test has not yet been executed.

\textbf{Coarse behavioral measure.} The 5-question sanity check produces
only two values (0.80 or 0.60). A continuous coherence metric would
strengthen the behavioral-change evidence.

% ============================================================================
\section{Conclusion}
\label{sec:conclusion}

We have introduced presence, a geometric metric for measuring identity
geometry during inference-time model interventions. The metric is
unsupervised, real-time, and validated by a same-layer positive control
that produces a clean dose-response from $1.000$ to $0.828$ under direct
identity-subspace injection.

Three results establish the contribution. First, V-cache injection at
input layers does not propagate to deep identity geometry: 168~trials
show presence at $0.978 \pm 0.002$ regardless of dose or strategy.
The same-layer positive control confirms this is genuine preservation,
not metric insensitivity: injection at L35 produces the clean
$1.000 \to 0.828$ dose-response. Second, the metric detects identity
perturbation when it is present (positive control) and correctly reports
preservation in the cross-layer experiment. Third, emotion vectors
share 22--25\% of their energy with the identity subspace at the
injection layers, yet identity at deep layers is unaffected.

The mechanism remains an open question. We hypothesize architectural
filtering by the hybrid model's linear-attention layers, but active
identity restoration cannot be ruled out with the current data. A
layer-matched smoke test designed to discriminate these mechanisms has
not yet been executed. This finding is specific to the GatedDeltaNet
architecture and may not generalize to dense transformers. For
inference-time correction safety, V-cache injection at input layers
preserves identity geometry \emph{in this architecture}; the exact
mechanism does not affect this practical finding.

The metric is available as open-source code and requires no specialized
hardware or training data.

% ============================================================================
\section*{Data and Code Availability}

All experiment scripts, preregistration documents, and raw results
(168-trial schema correction, 72-trial positive control, orthogonality
analysis, overnight validation) are available at
\url{https://github.com/Liberation-Labs-THCoalition/Project-Oracle}
(private repository; correction vectors and calibration tools withheld
under staged safety release due to dual-use risk; detection-side code
and experimental results available upon request to verified safety
researchers).
The model (Jackrong/Qwen3.5-27B-Claude-4.6-Opus-Reasoning-Distilled) and
SAE features (Qwen-Scope) are publicly available on HuggingFace.

% ============================================================================
\bibliographystyle{plainnat}
\bibliography{references}

\end{document}
